\chapter{Podsumowanie}

Celem pracy było zaprojektowanie, implementacja oraz przeprowadzenie analizy statystycznej
generatora bitów losowych opartego na oscylatorach pierścieniowych w strukturze układu
programowalnego FPGA. Cel ten został osiągnięty, a zakres wykonanych prac objął zarówno
warstwę sprzętową, jak i narzędzia programowe niezbędne do rejestracji i oceny
uzyskiwanych danych.

W części sprzętowej opracowano kompletny projekt w języku VHDL, obejmujący
parametryzowany moduł oscylatora pierścieniowego, moduł ekstrakcji entropii oparty na
detekcji zboczy w próbkowanych przebiegach oraz nadajnik interfejsu UART zrealizowany
w postaci automatu skończonego. Poszczególne bloki połączono mechanizmem potwierdzania
gotowości, co pozwoliło oddzielić warstwę generującą entropię od warstwy transmisyjnej
i uzależnić tempo wystawiania danych od przepustowości interfejsu komunikacyjnego.
Struktury oscylatorów objęto atrybutami zabezpieczającymi je przed usunięciem lub
uproszczeniem na etapie syntezy, a ich rozmieszczenie w obrębie układu określono ręcznie
ograniczeniami umieszczonymi w pliku ograniczeń projektu, co zapewniło powtarzalność
implementacji pomiędzy kolejnymi przebiegami narzędzia.

W części programowej przygotowano zestaw skryptów w języku Python, tworzących spójne
środowisko badawcze. Wydzielenie modułu parametrów konfiguracyjnych, oddzielenie etapu
akwizycji od etapów analizy oraz zapis serii pomiarowych do plików pozwoliły na
wielokrotną analizę tych samych danych przy różnych ustawieniach, bez konieczności
ponownego uruchamiania układu. Rozwiązanie to okazało się istotne dla wiarygodności
badań, ponieważ umożliwiło bezpośrednie porównanie strumienia surowego z poddanym
korekcji rozkładu w obrębie tej samej rejestracji.

Zasadniczym rezultatem pracy nie jest jednak sama implementacja, lecz wnioski wynikające
z porównania czterech konfiguracji generatora trzema niezależnymi metodami analizy
statystycznej. Badania wykazały, że poszczególne metody prowadzą do odmiennego
uszeregowania konfiguracji, a zabiegi projektowe poprawiające jedną właściwość strumienia
mogą jednocześnie pogarszać inną. Rozproszenie oscylatorów w strukturze układu
zredukowało odchylenie rozkładu bitów do poziomu naturalnych fluktuacji statystycznych,
pogarszając jednocześnie wynik badania autokorelacji o blisko dwa rzędy wielkości.
Praktycznym wnioskiem wynikającym z tej obserwacji jest konieczność oceny źródeł entropii
zestawem metod badających różne aspekty strumienia, ponieważ pozytywny wynik pojedynczego
testu, w szczególności powszechnie stosowanego testu monobitowego, nie stanowi podstawy
do stwierdzenia poprawności generatora.

Analiza pozwoliła również zlokalizować mechanizm odpowiedzialny za zasadniczą część
wykrytych odstępstw od losowości. Wykryte zależności koncentrowały się w zakresie
opóźnień odpowiadających liczbie kanałów wejściowych oraz długości bajtu, a ich charakter
zmieniał się wraz z relacjami częstotliwości pomiędzy oscylatorami, nie zaś wraz
z jakością samych źródeł szumu. Wskazuje to na warstwę serializacji, w której bity
zapisywane są do bufora w stałej kolejności numerów kanałów, jako na miejsce powstawania
struktury okresowej w strumieniu wyjściowym. Jest to wniosek o znaczeniu praktycznym,
ponieważ przesuwa uwagę projektanta z optymalizacji samych oscylatorów na sposób łączenia
danych pochodzących z wielu równoległych źródeł.

Praca ma określone ograniczenia, które warto wskazać. Przejście pomiędzy dwiema
konfiguracjami obejmowało jednoczesną zmianę dwóch parametrów, w związku z czym nie jest
możliwe rozdzielenie ich wpływu na obserwowaną poprawę. Długość serii pomiarowej,
wynikająca z przepustowości zastosowanego interfejsu komunikacyjnego, wyznacza granicę
czułości obu testów ilościowych i sprawia, że mapy dwuwymiarowe dla słów ośmiobitowych
zachowują wartość jedynie jakościową. Badania przeprowadzono ponadto w stałych warunkach
otoczenia, bez oceny wpływu temperatury oraz napięcia zasilania na stabilność źródła
entropii, co w przypadku generatorów sprzętowych stanowi istotny aspekt ich
charakterystyki. Nie zastosowano również standaryzowanych pakietów testów, których
wykorzystanie pozwoliłoby odnieść uzyskane wyniki do wymagań formalnych stawianych
źródłom losowości.

Naturalnym kierunkiem dalszych prac jest weryfikacja hipotezy dotyczącej warstwy
serializacji przez rozdzielenie zarejestrowanego strumienia na podstrumienie
odpowiadające poszczególnym kanałom oraz zbadanie każdego z nich osobno. Kolejnym jest
modyfikacja samego mechanizmu łączenia bitów, na przykład przez wprowadzenie zmiennej
kolejności zapisu kanałów do bufora, co powinno usunąć zależności okresowe u ich źródła
bez ponoszenia kosztu przepustowości właściwego dla korekcji von Neumanna. Wartościowe
byłoby także przeniesienie korekcji do struktury układu FPGA i umieszczenie jej
bezpośrednio przed nadajnikiem, co pozwoliłoby uzyskać strumień o poprawionych
właściwościach statystycznych bez przetwarzania danych po stronie komputera, oraz
zwiększenie przepustowości interfejsu komunikacyjnego, umożliwiające rejestrację serii
na tyle długich, by rozszerzyć zakres stosowalności zastosowanych metod na dłuższe słowa
bitowe.